Mamba-2 establishes theoretical equivalence between Transformer attention and state space models, promising a principled path for knowledge transfer from pretrained Transformers to efficient SSM architectures. We test whether this duality enables practical distillation by extracting SSM parameters directly from BERT attention weights. Our experiments reveal a surprising finding: while duality-derived parameters are numerically stable (0\% NaN/Inf, 0.11$\times$ magnitude ratio), they reconstruct attention outputs 2.13\% worse than random initialization---a consistent negative effect across all 12 BERT layers ($p < 0.0001$, Cohen's $d = -4.56$). We identify probable causes: the output-level Frobenius norm metric may not capture what duality preserves, and BERT's multi-head attention likely violates the theoretical SSD conditions. Our work provides the first empirical scope clarification for Mamba-2 duality: valid parameters do not imply preserved structure. Practitioners should test reconstruction quality, not just parameter validity, when developing attention-to-SSM conversion methods.
